\documentclass{ximera}

\input{../../../preamble.tex}


\definecolor{fvdnavy}{HTML}{071D43}
\definecolor{fvdcyan}{HTML}{20BFC6}
\definecolor{fvdcyanlight}{HTML}{EFFCFB}
\definecolor{fvdmagenta}{HTML}{F10D69}
\definecolor{fvdmagentalight}{HTML}{FFF1F7}
\definecolor{fvdtext}{HTML}{163044}





%=================================================
% Pedagogische omgevingen
% PDF: tcolorbox
% HTML: learning-box
%=================================================

\newenvironment{voorkennis}
{%
    \ifdefined\HCode
        \HCode{%
<section class="learning-box learning-box--prior">
<div class="learning-box__header">
<span class="learning-box__icon" aria-hidden="true"></span>
<span class="learning-box__title">Voorkennis</span>
</div>
<div class="learning-box__content">
        }%
        \begin{itemize}
    \else
       \begin{tcolorbox}[
    enhanced,
    breakable,
    colback=fvdcyanlight,
    colframe=fvdcyan,
    coltext=fvdtext,
    boxrule=0.5pt,
    leftrule=4pt,
    arc=4mm,
    outer arc=4mm,
    left=7mm,
    right=6mm,
    top=5mm,
    bottom=5mm,
    before skip=6mm,
    after skip=6mm,
    title=Voorkennis,
    coltitle=fvdcyan,
    fonttitle=\bfseries\Large,
    attach title to upper={\par\smallskip},
    boxed title style={
        empty,
        left=0mm,
        right=0mm,
        top=0mm,
        bottom=0mm
    },
    overlay={
        \draw[
            fvdcyan,
            line width=0.5pt
        ]
        ([xshift=42mm,yshift=-13mm]frame.north west)
        --
        ([xshift=-7mm,yshift=-13mm]frame.north east);
    }
]
        \begin{itemize}
    \fi
}
{%
    \end{itemize}
    \ifdefined\HCode
        \HCode{%
</div>
</section>
        }%
    \else
        \end{tcolorbox}
    \fi
}


\newenvironment{leerdoelen}
{%
    \ifdefined\HCode
        \HCode{%
<section class="learning-box learning-box--goals">
<div class="learning-box__header">
<span class="learning-box__icon" aria-hidden="true"></span>
<span class="learning-box__title">Leerdoelen</span>
</div>
<div class="learning-box__content">
        }%
        \begin{itemize}
    \else
\begin{tcolorbox}[
    enhanced,
    breakable,
    colback=fvdmagentalight,
    colframe=fvdmagenta,
    coltext=fvdtext,
    boxrule=0.5pt,
    leftrule=4pt,
    arc=4mm,
    outer arc=4mm,
    left=7mm,
    right=6mm,
    top=5mm,
    bottom=5mm,
    before skip=6mm,
    after skip=6mm,
    title=Leerdoelen,
    coltitle=fvdmagenta,
    fonttitle=\bfseries\Large,
    attach title to upper={\par\smallskip},
    boxed title style={
        empty,
        left=0mm,
        right=0mm,
        top=0mm,
        bottom=0mm
    },
    overlay={
        \draw[
            fvdmagenta,
            line width=0.5pt
        ]
        ([xshift=42mm,yshift=-13mm]frame.north west)
        --
        ([xshift=-7mm,yshift=-13mm]frame.north east);
    }
]
        \begin{itemize}
    \fi
}
{%
    \end{itemize}
    \ifdefined\HCode
        \HCode{%
</div>
</section>
        }%
    \else
        \end{tcolorbox}
    \fi
}


\addPrintStyle{../../..}

\title{Limieten en asymptotisch gedrag}
\author{Ingmar Herreman}

\begin{abstract}
Een functie hoeft niet overal een functiewaarde te hebben om toch
voorspelbaar gedrag te vertonen. Met limieten beschrijven we wat er
met de functiewaarden gebeurt wanneer $x$ een bepaalde waarde nadert
of onbeperkt groot of klein wordt.

In het vijfde jaar leerde je al limieten berekenen en asymptoten
bepalen bij veeltermfuncties, rationale en irrationale functies.
In dit hoofdstuk halen we de kern daarvan opnieuw naar boven.

De nadruk ligt niet op nieuwe rekentechnieken, maar op betekenis:
wat vertelt een limiet over een grafiek en hoe herken je daaruit
asymptotisch gedrag?
\end{abstract}

\begin{document}

% FVD_CHAPTER_DATA_START
% Automatisch gegenereerd uit index.tex.
% Niet handmatig aanpassen.
\fvdchapterdata{1}{Veranderingen in beeld}{4}
% FVD_CHAPTER_DATA_END









































































\maketitle


% ============================================================
% VOORKENNIS EN LEERDOELEN
% ============================================================

\begin{voorkennis}
    \item Je kent het domein van veelterm-, rationale en irrationale functies.
    \item Je kan eenvoudige limieten algebraïsch berekenen.
    \item Je kan rationale uitdrukkingen vereenvoudigen.
    \item Je kan bij wortelvormen werken met de toegevoegde factor.
    \item Je kan verticale, horizontale en schuine asymptoten bepalen.
    \item Je kan informatie aflezen uit een grafiek.
  \end{voorkennis}
  
  \begin{leerdoelen}
    \item Je kan in woorden uitleggen wat een limiet beschrijft.
    \item Je kan functiewaarde en limiet van elkaar onderscheiden.
    \item Je kan linker- en rechterlimieten interpreteren.
    \item Je kan limieten in een punt en op oneindig uit een grafiek aflezen.
    \item Je kan een oneindige limiet correct interpreteren.
    \item Je kan verticale, horizontale en schuine asymptoten herkennen.
    \item Je kan het verband tussen limieten en asymptotisch gedrag uitleggen.
    \item Je kan eerder geleerde limiettechnieken opnieuw gebruiken.
    \item Je kan vanuit limietinformatie uitspraken doen over het mogelijke verloop van een grafiek.
  \end{leerdoelen}


% ============================================================
% 1 WAAROM?
% ============================================================

\FVDsection{Waarom opnieuw limieten?}

In de vorige hoofdstukken onderzochten we hoe een functie
\emph{verandert}.

Met $f'$ en $f''$ kunnen we bijvoorbeeld stijgen, dalen, extrema en
kromming onderzoeken.

Maar daarmee kennen we nog niet het volledige gedrag van een functie.

Bij de functie

\[
f(x)=\frac{1}{x-2}
\]

is bijvoorbeeld vooral interessant wat er gebeurt wanneer

\[
x\to2
\]

of wanneer

\[
x\to+\infty.
\]

De functie is niet gedefinieerd voor $x=2$, maar haar gedrag in de
buurt van $2$ kunnen we wel heel precies beschrijven.

Daarvoor gebruiken we \textbf{limieten}.

\begin{opmerking}
Een limiet beantwoordt niet in de eerste plaats de vraag

\[
\text{``Wat is de functiewaarde?''}
\]

maar wel

\[
\boxed{
\text{``Waarheen gaan de functiewaarden wanneer }x
\text{ een bepaalde waarde nadert?''}
}
\]
\end{opmerking}


% ============================================================
% 2 KERNIDEE
% ============================================================

\FVDsection{Wat vertelt een limiet?}

Neem een functie $f$ en laat $x$ steeds dichter bij een getal $a$
komen.

Wanneer de functiewaarden daarbij steeds dichter bij een getal $L$
komen, schrijven we

\[
\boxed{
\lim_{x\to a}f(x)=L.
}
\]

We lezen:

\begin{center}
``de limiet van $f(x)$ voor $x$ gaande naar $a$ is $L$''.
\end{center}

Het getal $a$ vertelt dus \emph{waarheen $x$ gaat}.

Het getal $L$ vertelt \emph{waarheen $f(x)$ gaat}.


% ============================================================
% GRAFIEK 1
% ============================================================

\FVDsubsection{Een limiet gaat over gedrag in de buurt}

\begin{center}
\begin{tikzpicture}[x=1.15cm,y=0.9cm]

  % assen
  \draw[->,black!45] (-3.2,0) -- (3.5,0)
    node[right] {$x$};
  \draw[->,black!45] (0,-0.8) -- (0,4.4)
    node[above] {$y$};

  % grafiek y = x^2 met gat bij x=1
  \draw[fvdnavy,very thick,smooth,domain=-2:0.94,samples=80]
    plot (\x,{0.7*(\x)^2+1});

  \draw[fvdnavy,very thick,smooth,domain=1.06:2.1,samples=80]
    plot (\x,{0.7*(\x)^2+1});

  % open punt
  \draw[fvdmagenta,very thick,fill=white]
    (1,1.7) circle (2.5pt);

  % aparte functiewaarde
  \fill[fvdmagenta]
    (1,3.3) circle (2.5pt);

  % hulplijnen
  \draw[dashed,black!45]
    (1,0) -- (1,1.7);

  \draw[dashed,black!45]
    (0,1.7) -- (1,1.7);

  \node[below] at (1,0) {$a$};
  \node[left] at (0,1.7) {$L$};

  \node[right] at (1.1,3.3) {$f(a)$};

\end{tikzpicture}
\end{center}

In deze grafiek geldt:

\[
f(a)\neq L,
\]

maar toch:

\[
\boxed{
\lim_{x\to a}f(x)=L.
}
\]

De waarden van $f(x)$ naderen immers $L$ wanneer $x$ naar $a$
nadert.

\begin{eigenschap}[Functiewaarde en limiet]

De uitspraken

\[
f(a)=L
\]

en

\[
\lim_{x\to a}f(x)=L
\]

zeggen niet hetzelfde.

De eerste gaat over de waarde \emph{in} $a$.

De tweede gaat over het gedrag van de functie \emph{in de buurt van}
$a$.

\end{eigenschap}


% ============================================================
% OEFENING
% ============================================================

\begin{oefening}[Functiewaarde of limiet?]{\oefB\ESS}

Uit een grafiek lezen we:

\[
f(2)=5
\]

en

\[
\lim_{x\to2}f(x)=3.
\]

Welke uitspraak is correct?

\begin{multipleChoice}
\choice{Dit is onmogelijk: functiewaarde en limiet moeten gelijk zijn.}
\choice{De grafiek gaat noodzakelijk door $(2,3)$.}
\choice[correct]{In $x=2$ is de functiewaarde $5$, terwijl de functiewaarden in de buurt van $2$ naar $3$ naderen.}
\choice{De functie is niet gedefinieerd in $x=2$.}
\end{multipleChoice}

\end{oefening}


% ============================================================
% 3 LINKER EN RECHTERLIMIET
% ============================================================

\FVDsection{Van links of van rechts naderen}

We kunnen een waarde $a$ op twee manieren naderen.

Van links:

\[
x\to a^-.
\]

Van rechts:

\[
x\to a^+.
\]

Daarom onderscheiden we:

\[
\boxed{
\lim_{x\to a^-}f(x)
}
\]

en

\[
\boxed{
\lim_{x\to a^+}f(x).
}
\]

Dit noemen we respectievelijk de \textbf{linkerlimiet} en de
\textbf{rechterlimiet}.

\begin{eigenschap}[Wanneer bestaat de tweezijdige limiet?]

Als beide eenzijdige limieten bestaan en gelijk zijn, dan bestaat de
limiet:

\[
\boxed{
\lim_{x\to a}f(x)=L
}
\]

precies wanneer

\[
\lim_{x\to a^-}f(x)=L
\]

en

\[
\lim_{x\to a^+}f(x)=L.
\]

Zijn linker- en rechterlimiet verschillend, dan bestaat de
tweezijdige limiet niet.

\end{eigenschap}


% ============================================================
% OEFENING
% ============================================================

\begin{oefening}[Links en rechts]{\oefB\ESS}

Voor een functie $f$ geldt:

\[
\lim_{x\to3^-}f(x)=2
\]

en

\[
\lim_{x\to3^+}f(x)=5.
\]

Welke uitspraak is correct?

\begin{multipleChoice}
\choice{$\displaystyle\lim_{x\to3}f(x)=2$}
\choice{$\displaystyle\lim_{x\to3}f(x)=5$}
\choice{$\displaystyle\lim_{x\to3}f(x)=\frac72$}
\choice[correct]{De tweezijdige limiet voor $x\to3$ bestaat niet.}
\end{multipleChoice}

\end{oefening}


% ============================================================
% 4 ONEINDIGE LIMIET
% ============================================================

\FVDsection{Wanneer functiewaarden onbegrensd worden}

Een limiet hoeft geen reëel getal te zijn.

Bijvoorbeeld:

\[
\lim_{x\to2^+}\frac{1}{x-2}=+\infty.
\]

Dat betekent:

wanneer $x$ langs rechts steeds dichter bij $2$ komt, worden de
functiewaarden positief en onbeperkt groot.

Langs links geldt:

\[
\lim_{x\to2^-}\frac{1}{x-2}=-\infty.
\]


% ============================================================
% GRAFIEK 2
% ============================================================

\begin{center}
\begin{tikzpicture}[x=1.15cm,y=0.8cm]

  \draw[->,black!45] (-2.5,0) -- (4.8,0)
    node[right] {$x$};
  \draw[->,black!45] (0,-4.2) -- (0,4.2)
    node[above] {$y$};

  % verticale asymptoot x=2
  \draw[dashed,draw=fvdmagenta,thick]
    (2,-4.0) -- (2,4.0);

  \node[below] at (2,0) {$2$};

  % 1/(x-2)
  \draw[fvdnavy,very thick,smooth,domain=-2.2:1.75,samples=100]
    plot (\x,{1/(\x-2)});

  \draw[fvdnavy,very thick,smooth,domain=2.25:4.6,samples=100]
    plot (\x,{1/(\x-2)});

\end{tikzpicture}
\end{center}

De rechte

\[
x=2
\]

beschrijft hier een verticale grens die de grafiek steeds dichter
nadert.

Dat brengt ons bij asymptoten.


% ============================================================
% 5 VERTICALE ASYMPTOTEN
% ============================================================

\FVDsection{Verticale asymptoten}

\begin{definitie}[Verticale asymptoot]

De rechte

\[
\boxed{x=a}
\]

is een \textbf{verticale asymptoot} van de grafiek van $f$ wanneer
minstens één van de relevante eenzijdige limieten bij $a$ oneindig
is.

Bijvoorbeeld:

\[
\lim_{x\to a^-}f(x)=+\infty,
\qquad
\lim_{x\to a^+}f(x)=-\infty.
\]

\end{definitie}

Let op: het teken links en rechts hoeft niet hetzelfde te zijn.

Bij

\[
f(x)=\frac{1}{(x-2)^2}
\]

geldt bijvoorbeeld:

\[
\lim_{x\to2^-}f(x)=+\infty
\]

en

\[
\lim_{x\to2^+}f(x)=+\infty.
\]

Ook daar is

\[
x=2
\]

een verticale asymptoot.


% ============================================================
% OEFENING
% ============================================================

\begin{oefening}[Van limieten naar verticale asymptoot]{\oefB\ESS}

Gegeven:

\[
\lim_{x\to-1^-}f(x)=-\infty
\]

en

\[
\lim_{x\to-1^+}f(x)=+\infty.
\]

Welke conclusie kun je trekken?

\begin{multipleChoice}
\choice{$y=-1$ is een horizontale asymptoot.}
\choice[correct]{$x=-1$ is een verticale asymptoot.}
\choice{$(-1,0)$ is noodzakelijk een nulpunt.}
\choice{$f(-1)=0$.}
\end{multipleChoice}

\end{oefening}


% ============================================================
% 6 GEDRAG OP ONEINDIG
% ============================================================

\FVDsection{Wat gebeurt er ver naar links en rechts?}

Limieten kunnen ook beschrijven wat er gebeurt wanneer $x$
onbegrensd groot of klein wordt.

We schrijven bijvoorbeeld

\[
\lim_{x\to+\infty}f(x)=3.
\]

Dit betekent dat de functiewaarden steeds dichter bij $3$ komen
wanneer we op de grafiek steeds verder naar rechts gaan.

Analoog beschrijft

\[
\lim_{x\to-\infty}f(x)
\]

het gedrag ver naar links.

\begin{opmerking}
De symbolen

\[
+\infty
\qquad\text{en}\qquad
-\infty
\]

zijn hier geen gewone getallen.

De notatie $x\to+\infty$ betekent dat $x$ onbeperkt toeneemt.
\end{opmerking}


% ============================================================
% 7 HORIZONTALE ASYMPTOTEN
% ============================================================

\FVDsection{Horizontale asymptoten}

\begin{definitie}[Horizontale asymptoot]

Als

\[
\lim_{x\to+\infty}f(x)=b
\]

of

\[
\lim_{x\to-\infty}f(x)=b,
\]

dan is

\[
\boxed{y=b}
\]

een \textbf{horizontale asymptoot} van de grafiek van $f$ aan de
overeenkomstige zijde.

\end{definitie}


% ============================================================
% GRAFIEK 3
% ============================================================

\begin{center}
\begin{tikzpicture}[x=1.0cm,y=0.8cm]

  \draw[->,black!45] (-4.5,0) -- (5,0)
    node[right] {$x$};
  \draw[->,black!45] (0,-1) -- (0,4.7)
    node[above] {$y$};

  % horizontale asymptoot
  \draw[dashed,fvdmagenta,thick]
    (-4.3,2) -- (4.8,2);

  \node[left] at (0,2) {$2$};

  % voorbeeldcurve y=2+1/(x^2+1)
  \draw[fvdnavy,very thick,smooth,domain=-4.2:4.6,samples=120]
    plot (\x,{2+1/(\x*\x+1)});

\end{tikzpicture}
\end{center}

Voor deze grafiek zien we:

\[
\lim_{x\to-\infty}f(x)=2
\]

en

\[
\lim_{x\to+\infty}f(x)=2.
\]

Dus:

\[
\boxed{y=2}
\]

is een horizontale asymptoot.


% ============================================================
% BELANGRIJKE DENKFOUT
% ============================================================

\FVDsubsection{Mag een grafiek een horizontale asymptoot snijden?}

Ja.

Een horizontale asymptoot beschrijft het gedrag voor

\[
x\to+\infty
\]

of

\[
x\to-\infty.
\]

Ze zegt niet dat de grafiek de rechte nergens mag snijden.

\begin{opmerking}
Een asymptoot is geen ``muur'' waar een grafiek niet doorheen mag.

De limiet beschrijft hoe de afstand tussen de grafiek en de
asymptoot zich gedraagt in de relevante richting.
\end{opmerking}


% ============================================================
% 8 SCHUINE ASYMPTOTEN
% ============================================================

\FVDsection{Schuine asymptoten}

Een functie hoeft voor grote $|x|$ niet naar een constante waarde te
naderen.

Ze kan ook steeds dichter bij een schuine rechte komen.

\begin{definitie}[Schuine asymptoot]

De rechte

\[
\boxed{y=ax+b}
\qquad (a\neq0)
\]

is een schuine asymptoot wanneer

\[
\boxed{
\lim_{x\to\pm\infty}
\bigl(f(x)-(ax+b)\bigr)=0
}
\]

voor de onderzochte richting.

\end{definitie}

Het verschil

\[
f(x)-(ax+b)
\]

is de verticale afstand met teken tussen de functiewaarde en de
waarde op de rechte.

Als dat verschil naar nul gaat, komen grafiek en rechte steeds
dichter bij elkaar.


% ============================================================
% GRAFIEK 4
% ============================================================

\begin{center}
\begin{tikzpicture}[x=0.9cm,y=0.7cm]

  \draw[->,black!45] (-4.8,0) -- (5.2,0)
    node[right] {$x$};
  \draw[->,black!45] (0,-3.5) -- (0,5.2)
    node[above] {$y$};

  % asymptoot y = 0.6x + 1
  \draw[dashed,fvdmagenta,thick,domain=-4.5:4.8]
    plot (\x,{0.6*\x+1});

  \node[fvdmagenta] at (3.5,3.8) {$y=ax+b$};

  % curve die asymptoot nadert
  \draw[fvdnavy,very thick,smooth,domain=-4.4:4.8,samples=120]
    plot (\x,{0.6*\x+1+1/(\x*\x+1)});

\end{tikzpicture}
\end{center}


% ============================================================
% OEFENING ASYMPTOTEN
% ============================================================

\begin{oefening}[Welke asymptoot?]{\oefB\ESS}

Koppel elke limietuitspraak aan de juiste conclusie.

\begin{question}

\[
\lim_{x\to4^+}f(x)=+\infty.
\]

Welke asymptoot kan hieruit volgen?

\begin{multipleChoice}
\choice{$y=4$}
\choice[correct]{$x=4$}
\choice{$y=x+4$}
\end{multipleChoice}

\end{question}

\begin{question}

\[
\lim_{x\to+\infty}f(x)=-3.
\]

Welke asymptoot volgt hieruit?

\begin{multipleChoice}
\choice{$x=-3$}
\choice[correct]{$y=-3$}
\choice{$y=x-3$}
\end{multipleChoice}

\end{question}

\begin{question}

\[
\lim_{x\to+\infty}
\bigl(f(x)-(2x-1)\bigr)=0.
\]

Welke asymptoot volgt hieruit?

\begin{multipleChoice}
\choice{$x=2$}
\choice{$y=-1$}
\choice[correct]{$y=2x-1$}
\end{multipleChoice}

\end{question}

\end{oefening}


% ============================================================
% 9 DE DRIE SOORTEN SAMEN
% ============================================================

\FVDsection{Drie soorten asymptotisch gedrag}

De drie situaties kunnen we nu naast elkaar plaatsen.

\[
\boxed{
\begin{array}{c|c}
\text{limietgedrag}
&
\text{asymptoot}
\\ \hline
f(x)\to\pm\infty
\text{ voor }x\to a^\pm
&
x=a
\\[2mm]
f(x)\to b
\text{ voor }x\to\pm\infty
&
y=b
\\[2mm]
f(x)-(ax+b)\to0
\text{ voor }x\to\pm\infty
&
y=ax+b
\end{array}
}
\]

\begin{opmerking}
Probeer deze verbanden niet als drie losse recepten te onthouden.

Vraag telkens:

\[
\boxed{\text{Wat doet de grafiek?}}
\]

Bij een verticale asymptoot wordt $x$ begrensd terwijl $f(x)$
onbegrensd wordt.

Bij een horizontale of schuine asymptoot gaat $x$ juist naar
$\pm\infty$ en nadert de grafiek steeds dichter tot een rechte.
\end{opmerking}


% ============================================================
% 10 REKENTECHNIEKEN
% ============================================================

\FVDsection{Kun je de limieten nog berekenen?}

In het vijfde jaar heb je verschillende technieken geleerd om
limieten algebraïsch te bepalen.

We bouwen die technieken hier niet opnieuw volledig op.

Je moet wel opnieuw herkennen welke aanpak geschikt is.

\begin{center}
\[
\boxed{
\text{eerst analyseren}
\quad\longrightarrow\quad
\text{dan pas rekenen}
}
\]
\end{center}

Een nuttige eerste stap is altijd:

\begin{enumerate}
  \item naar welke waarde gaat $x$?
  \item wat verwacht je dat teller en noemer doen?
  \item ontstaat een gewone waarde, oneindig gedrag of een
        onbepaalde vorm?
  \item welke eerder geleerde techniek kan die vorm vereenvoudigen?
\end{enumerate}


% ============================================================
% VEELTERMFUNCTIE
% ============================================================

\FVDsubsection{Veeltermfuncties}

Voor gedrag op oneindig bepaalt de term van de hoogste graad het
uiteindelijke gedrag.

Bijvoorbeeld:

\[
f(x)=-2x^5+3x^2-7.
\]

Voor grote positieve $x$ overheerst

\[
-2x^5.
\]

Daarom:

\[
\boxed{
\lim_{x\to+\infty}f(x)=-\infty.
}
\]

Voor $x\to-\infty$ geldt:

\[
-2x^5\to+\infty,
\]

dus:

\[
\boxed{
\lim_{x\to-\infty}f(x)=+\infty.
}
\]


% ============================================================
% RATIONALE FUNCTIE
% ============================================================

\FVDsubsection{Rationale functies}

Neem

\[
f(x)=\frac{2x^2+3x-1}{x^2-4}.
\]

Voor $|x|$ groot delen we teller en noemer bijvoorbeeld door $x^2$:

\[
f(x)
=
\frac{
2+\frac3x-\frac1{x^2}
}{
1-\frac4{x^2}
}.
\]

Voor $x\to\pm\infty$ verdwijnen de termen met $\frac1x$ en
$\frac1{x^2}$.

Dus:

\[
\boxed{
\lim_{x\to\pm\infty}f(x)=2.
}
\]

Daaruit volgt de horizontale asymptoot

\[
\boxed{y=2.}
\]

Voor verticale asymptoten onderzoeken we onder meer waar de noemer
nul wordt en wat de functie daar langs beide zijden doet.


% ============================================================
% IRRATIONALE FUNCTIE
% ============================================================

\FVDsubsection{Irrationale functies}

Bij wortelvormen kan de toegevoegde factor opnieuw nuttig zijn.

Neem bijvoorbeeld

\[
f(x)=\sqrt{x^2+4x}-x,
\qquad x\to+\infty.
\]

Rechtstreeks invullen suggereert de onbepaalde vorm

\[
\infty-\infty.
\]

We vermenigvuldigen met de toegevoegde factor:

\[
\sqrt{x^2+4x}-x
=
\frac{
(\sqrt{x^2+4x}-x)(\sqrt{x^2+4x}+x)
}{
\sqrt{x^2+4x}+x
}.
\]

Dus:

\[
f(x)
=
\frac{4x}
{\sqrt{x^2+4x}+x}.
\]

Voor $x>0$ kunnen we $x$ afzonderen:

\[
f(x)
=
\frac{4}
{\sqrt{1+\frac4x}+1}.
\]

Daarom:

\[
\boxed{
\lim_{x\to+\infty}f(x)=2.
}
\]

De grafiek heeft voor $x\to+\infty$ dus horizontale asymptoot

\[
\boxed{y=2.}
\]


% ============================================================
% DIAGNOSTISCHE OEFENING
% ============================================================

\begin{oefening}[Parate rekentechniek]{\oefB\ESS}

Bereken zonder dat een methode wordt aangegeven.

\begin{question}

\[
\lim_{x\to+\infty}
\left(-3x^4+2x^2-7\right).
\]

\end{question}

\begin{question}

\[
\lim_{x\to+\infty}
\frac{5x^2-x+1}{2x^2+3}.
\]

\end{question}

\begin{question}

\[
\lim_{x\to+\infty}
\left(\sqrt{x^2+6x}-x\right).
\]

\end{question}

Noteer bij elke limiet niet alleen het resultaat, maar ook welke
eigenschap of techniek je gebruikt.

\end{oefening}


% ============================================================
% 11 GRAFIEK -> LIMIET
% ============================================================

\FVDsection{Van grafiek naar limiet}

Limieten moeten niet altijd algebraïsch worden berekend.

Uit een grafiek kun je vaak rechtstreeks aflezen:

\begin{itemize}
  \item welke waarde $f(x)$ nadert voor $x\to a^-$;
  \item welke waarde $f(x)$ nadert voor $x\to a^+$;
  \item of de functiewaarden onbegrensd worden;
  \item hoe de functie zich gedraagt voor $x\to\pm\infty$;
  \item welke asymptoten zichtbaar zijn.
\end{itemize}

Belangrijk is dat je de grafiek \emph{gericht} leest.

Volg de grafiek in de richting waarin $x$ nadert.


% ============================================================
% 12 LIMIET -> GRAFIEK
% ============================================================

\FVDsection{Van limietinformatie naar een grafiek}

De omgekeerde richting is minstens even belangrijk.

Stel dat een functie $f$ voldoet aan:

\[
\lim_{x\to-2^-}f(x)=+\infty,
\]

\[
\lim_{x\to-2^+}f(x)=-\infty,
\]

\[
\lim_{x\to-\infty}f(x)=1
\]

en

\[
\lim_{x\to+\infty}f(x)=3.
\]

Zonder het voorschrift van $f$ te kennen, weten we al heel wat.

Er is een verticale asymptoot

\[
x=-2.
\]

Links nadert de grafiek voor $x\to-\infty$ de rechte

\[
y=1.
\]

Rechts nadert de grafiek voor $x\to+\infty$ de rechte

\[
y=3.
\]

De limieten leggen dus belangrijke randvoorwaarden op aan de vorm
van de grafiek.


% ============================================================
% GRAFIEKCONSTRUCTIE
% ============================================================

\begin{oefening}[Schets uit limietgegevens]{\oefU\ESS}

Van een functie $f$ is gegeven:

\[
\lim_{x\to1^-}f(x)=-\infty,
\]

\[
\lim_{x\to1^+}f(x)=+\infty,
\]

\[
\lim_{x\to-\infty}f(x)=2
\]

en

\[
\lim_{x\to+\infty}
\bigl(f(x)-(x-1)\bigr)=0.
\]

\begin{question}
Bepaal alle asymptoten die uit deze informatie volgen.
\end{question}

\begin{question}
Geef bij elke asymptoot aan aan welke zijde van de grafiek de
informatie betrekking heeft.
\end{question}

\begin{question}
Maak een mogelijke schets van de grafiek.

Je schets hoeft niet uniek te zijn, maar moet met alle gegeven
limieten overeenstemmen.
\end{question}

\begin{question}
Leg uit waarom verschillende grafieken aan dezelfde vier
limietvoorwaarden kunnen voldoen.
\end{question}

\end{oefening}


% ============================================================
% 13 EEN ASYMPTOOT IS GEEN VERBODSLIJN
% ============================================================

\FVDsection{Wat een asymptoot niet zegt}

Een asymptoot vertelt hoe een grafiek zich in een bepaalde
limietsituatie gedraagt.

Daaruit mag je niet méér besluiten dan de limiet werkelijk zegt.

\begin{opmerking}

\textbf{Een verticale asymptoot}

\[
x=a
\]

betekent niet automatisch dat de linker- en rechterlimiet hetzelfde
teken hebben.

\medskip

\textbf{Een horizontale asymptoot}

\[
y=b
\]

betekent niet dat de grafiek de rechte nooit kan snijden.

\medskip

\textbf{Een schuine asymptoot}

\[
y=ax+b
\]

betekent niet dat de grafiek voor eindige $x$ exact op die rechte
ligt.

\medskip

Een asymptoot beschrijft \textbf{naderingsgedrag}.

\end{opmerking}


% ============================================================
% 14 COMBINATIEOEFENING
% ============================================================

\begin{oefening}[Wat kun je werkelijk besluiten?]{\oefU\ESS}

Een functie $f$ voldoet aan

\[
\lim_{x\to+\infty}f(x)=4.
\]

Beoordeel de volgende uitspraken.

\begin{question}
Voor voldoende grote $x$ is $f(x)=4$.

\begin{multipleChoice}
\choice{Altijd waar.}
\choice[correct]{Niet noodzakelijk waar.}
\end{multipleChoice}
\end{question}

\begin{question}
De rechte $y=4$ is een horizontale asymptoot voor
$x\to+\infty$.

\begin{multipleChoice}
\choice[correct]{Waar.}
\choice{Niet waar.}
\end{multipleChoice}
\end{question}

\begin{question}
De grafiek kan de rechte $y=4$ nergens snijden.

\begin{multipleChoice}
\choice{Waar.}
\choice[correct]{Niet noodzakelijk waar.}
\end{multipleChoice}
\end{question}

\begin{question}
Uit deze limiet alleen kunnen we besluiten dat $f$ stijgend is voor
grote $x$.

\begin{multipleChoice}
\choice{Waar.}
\choice[correct]{Niet noodzakelijk waar.}
\end{multipleChoice}
\end{question}

\end{oefening}


% ============================================================
% 15 STEM
% ============================================================

\FVDsection{Asymptotisch gedrag in een model}

Asymptoten komen niet alleen voor in abstracte functievoorschriften.

Ze kunnen in een model een grenswaarde voorstellen die steeds dichter
wordt benaderd.

Stel bijvoorbeeld dat de temperatuur van een voorwerp wordt
beschreven door $T(t)$ en dat

\[
\lim_{t\to+\infty}T(t)=20.
\]

Dan betekent dit dat de temperatuur bij langdurig verloop steeds
dichter bij

\[
20\,^\circ\mathrm{C}
\]

komt.

De horizontale asymptoot

\[
T=20
\]

heeft dan een concrete betekenis: ze stelt de
\textbf{evenwichtstemperatuur} van het model voor.

\begin{opmerking}
Bij een toepassing moet je een limiet daarom altijd terug vertalen
naar de oorspronkelijke grootheden en eenheden.

Niet alleen

\[
\lim_{t\to+\infty}T(t)=20,
\]

maar bijvoorbeeld:

\begin{quote}
Op lange termijn nadert de temperatuur volgens het model
$20\,^\circ\mathrm{C}$.
\end{quote}
\end{opmerking}


% ============================================================
% 16 DENKVRAAG
% ============================================================

\begin{oefening}[Een grenswaarde interpreteren]{\oefU\ESS}

De concentratie $C(t)$ van een stof in een reactor wordt uitgedrukt
in $\mathrm{mg/L}$.

Er geldt:

\[
\lim_{t\to+\infty}C(t)=8.
\]

\begin{question}
Geef een correcte interpretatie van deze limiet in de context.
\end{question}

\begin{question}
Welke horizontale asymptoot hoort bij het model?
\end{question}

\begin{question}
Kun je uit deze informatie alleen besluiten dat de concentratie
altijd kleiner is dan $8\,\mathrm{mg/L}$?

Verantwoord.
\end{question}

\end{oefening}


% ============================================================
% 17 VEELGEMAAKTE DENKFOUTEN
% ============================================================

\FVDsection{Veelgemaakte denkfouten}

\begin{opmerking}

\textbf{Denkfout 1}

\[
\lim_{x\to a}f(x)=L
\]

betekent niet automatisch

\[
f(a)=L.
\]

\medskip

\textbf{Denkfout 2}

Als

\[
f(a)
\]

niet bestaat, betekent dat niet automatisch dat

\[
\lim_{x\to a}f(x)
\]

niet bestaat.

\medskip

\textbf{Denkfout 3}

Een nulwaarde van de noemer betekent niet zonder verder onderzoek
automatisch dat er een verticale asymptoot is.

Onderzoek het werkelijke limietgedrag.

\medskip

\textbf{Denkfout 4}

Een horizontale asymptoot mag door de grafiek worden gesneden.

\medskip

\textbf{Denkfout 5}

Bij

\[
x\to-\infty
\]

mag je tekens niet gedachteloos behandelen alsof $x$ positief is.

Dat is vooral belangrijk bij uitdrukkingen met wortels, want

\[
\sqrt{x^2}=|x|.
\]

\medskip

\textbf{Denkfout 6}

Een asymptoot is geen afzonderlijk rekenobject.

Ze volgt uit het limietgedrag van de functie.

\end{opmerking}


% ============================================================
% 18 STERKE SYNTHESE
% ============================================================

\FVDsection{Kun je alles verbinden?}

% ============================================================
% SYNTHESEGRAFIEK — LIMIETEN EN ASYMPTOTEN
% ============================================================

\FVDsection{Eén grafiek, meerdere limieten}

Een grafiek kan tegelijk verschillende vormen van asymptotisch gedrag
vertonen.

Beschouw de functie

\[
f(x)=x+1+\frac{1}{x-1}.
\]

De grafiek heeft een verticale asymptoot en een schuine asymptoot.

\begin{center}
\begin{tikzpicture}[x=0.9cm,y=0.7cm]

  % ----------------------------------------------------------
  % Assen
  % ----------------------------------------------------------

  \draw[->,black!45]
    (-5,0) -- (5.5,0)
    node[right] {$x$};

  \draw[->,black!45]
    (0,-4.5) -- (0,6)
    node[above] {$y$};


  % ----------------------------------------------------------
  % Verticale asymptoot x = 1
  % ----------------------------------------------------------

  \draw[
    dashed,
    fvdmagenta,
    thick
  ]
    (1,-4.2) -- (1,5.7);

  \node[
    fvdmagenta,
    below
  ]
    at (1,0)
    {$1$};


  % ----------------------------------------------------------
  % Schuine asymptoot y = x + 1
  % ----------------------------------------------------------

  \draw[
    dashed,
    fvdmagenta,
    thick,
    domain=-4.7:4.8
  ]
    plot (\x,{\x+1});

  \node[
    fvdmagenta,
    rotate=34
  ]
    at (3.7,5.0)
    {$y=x+1$};


  % ----------------------------------------------------------
  % Grafiek links van x=1
  % f(x)=x+1+1/(x-1)
  % ----------------------------------------------------------

  \draw[
    fvdnavy,
    very thick,
    smooth,
    domain=-4.7:0.78,
    samples=160
  ]
    plot (\x,{\x+1+1/(\x-1)});


  % ----------------------------------------------------------
  % Grafiek rechts van x=1
  % ----------------------------------------------------------

  \draw[
    fvdnavy,
    very thick,
    smooth,
    domain=1.22:4.8,
    samples=160
  ]
    plot (\x,{\x+1+1/(\x-1)});


  % ----------------------------------------------------------
  % Labels
  % ----------------------------------------------------------

  \node[
    fvdnavy
  ]
    at (-2.8,-2.3)
    {$f$};

\end{tikzpicture}
\end{center}


\begin{oefening}[Lees het asymptotisch gedrag]{\oefU\ESS}

Gebruik uitsluitend de grafiek.

\begin{question}

Wat is

\[
\lim_{x\to1^-}f(x)?
\]

\[
\answer{-\infty}
\]

\end{question}

\begin{question}

Wat is

\[
\lim_{x\to1^+}f(x)?
\]

\[
\answer{+\infty}
\]

\end{question}

\begin{question}

Welke verticale asymptoot volgt hieruit?

\[
\boxed{x=\answer{1}}
\]

\end{question}

\begin{question}

Wat gebeurt er met het verschil

\[
f(x)-(x+1)
\]

wanneer

\[
x\to+\infty?
\]

\[
\answer{0}
\]

\end{question}

\begin{question}

En wanneer

\[
x\to-\infty?
\]

\[
\answer{0}
\]

\end{question}

\begin{question}

Welke schuine asymptoot volgt daaruit?

\[
\boxed{y=\answer{x+1}}
\]

\end{question}

\end{oefening}

\begin{oefening}[Van formule naar volledig asymptotisch beeld]{\oefU\ESS}

Gegeven:

\[
f(x)=\frac{2x^2-x+3}{x-1}.
\]

\begin{question}
Bepaal het domein.
\end{question}

\begin{question}
Onderzoek het gedrag voor

\[
x\to1^-
\qquad\text{en}\qquad
x\to1^+.
\]

Bepaal daaruit een eventuele verticale asymptoot.
\end{question}

\begin{question}
Voer de veeltermdeling uit en schrijf $f$ in de vorm

\[
f(x)=ax+b+\frac{r}{x-1}.
\]

\end{question}

\begin{question}
Gebruik die vorm om het gedrag voor

\[
x\to\pm\infty
\]

te analyseren.
\end{question}

\begin{question}
Bepaal de schuine asymptoot.
\end{question}

\begin{question}
Maak op basis van je resultaten een kwalitatieve schets van de
grafiek.

Duid de asymptoten duidelijk aan.
\end{question}

\end{oefening}

\begin{opmerking}

    Uit één grafiek kunnen dus verschillende soorten limietinformatie
    worden afgeleid.
    
    Bij
    
    \[
    x\to1^\pm
    \]
    
    onderzoeken we lokaal gedrag in de buurt van een bepaalde
    $x$-waarde.
    
    Bij
    
    \[
    x\to\pm\infty
    \]
    
    onderzoeken we het gedrag van de grafiek ver naar links en rechts.
    
    Dezelfde grafiek kan dus tegelijk een verticale én een schuine
    asymptoot hebben.
    
    \end{opmerking}

% ============================================================
% 19 GEVORDERD
% ============================================================

\begin{oefening}[Kan dezelfde functie dit allemaal?]{\oefG}

Onderzoek of er een functie kan bestaan waarvoor tegelijk geldt:

\[
\lim_{x\to2^-}f(x)=+\infty,
\]

\[
\lim_{x\to2^+}f(x)=-\infty,
\]

\[
\lim_{x\to-\infty}f(x)=1
\]

en

\[
\lim_{x\to+\infty}\bigl(f(x)-(2x+3)\bigr)=0.
\]

Als je denkt dat dit mogelijk is, maak dan een kwalitatieve schets.

Als je denkt dat dit onmogelijk is, leg dan uit welke voorwaarden
elkaar tegenspreken.

\end{oefening}


% ============================================================
% 20 SAMENVATTING
% ============================================================

\FVDsection{Samengevat}

Een limiet beschrijft het gedrag van functiewaarden wanneer $x$ een
bepaalde waarde nadert:

\[
\boxed{
\lim_{x\to a}f(x)=L.
}
\]

Functiewaarde en limiet zijn verschillende begrippen:

\[
f(a)
\qquad\text{tegenover}\qquad
\lim_{x\to a}f(x).
\]

Voor een tweezijdige limiet moeten linker- en rechterlimiet
overeenkomen.

Een verticale asymptoot heeft vergelijking

\[
\boxed{x=a}
\]

en hangt samen met oneindig gedrag voor $x\to a^\pm$.

Een horizontale asymptoot heeft vergelijking

\[
\boxed{y=b}
\]

wanneer

\[
\lim_{x\to\pm\infty}f(x)=b.
\]

Een schuine asymptoot heeft vergelijking

\[
\boxed{y=ax+b}
\]

wanneer

\[
\lim_{x\to\pm\infty}
\bigl(f(x)-(ax+b)\bigr)=0.
\]

De centrale gedachte is:

\[
\boxed{
\text{limiet}
\quad\longleftrightarrow\quad
\text{gedrag van de grafiek}
}
\]

Een goede analyse gaat daarom in beide richtingen:

\[
\text{formule}
\longrightarrow
\text{limiet}
\longrightarrow
\text{grafiek}
\]

maar ook

\[
\text{grafiek}
\longrightarrow
\text{limietinformatie}.
\]


% ============================================================
% 21 MINIMUMTRAJECT
% ============================================================

\FVDsection{Wat moet je zeker beheersen?}

Na deze opfrissing moet je opnieuw vlot kunnen:

\begin{itemize}

  \item uitleggen wat een limiet betekent;

  \item functiewaarde en limiet onderscheiden;

  \item linker- en rechterlimieten lezen en interpreteren;

  \item limieten uit een grafiek aflezen;

  \item oneindige limieten interpreteren;

  \item het verband leggen tussen limieten en verticale,
        horizontale en schuine asymptoten;

  \item eenvoudige limieten van veelterm-, rationale en irrationale
        functies opnieuw berekenen;

  \item een geschikte rekentechniek zelfstandig herkennen;

  \item uit limietgegevens eigenschappen van een mogelijke grafiek
        afleiden;

  \item asymptotisch gedrag in een eenvoudige wetenschappelijke
        context interpreteren.

\end{itemize}

\begin{opmerking}
De oefeningen met de diamant vormen het minimumtraject.

Concreet:

\begin{itemize}
  \item \textit{Functiewaarde of limiet?};
  \item \textit{Links en rechts};
  \item \textit{Van limieten naar verticale asymptoot};
  \item \textit{Welke asymptoot?};
  \item \textit{Parate rekentechniek};
  \item \textit{Schets uit limietgegevens};
  \item \textit{Wat kun je werkelijk besluiten?};
  \item \textit{Een grenswaarde interpreteren};
  \item \textit{Van formule naar volledig asymptotisch beeld}.
\end{itemize}

Merk op dat meerdere essentiële oefeningen als \oefU{} zijn
geclassificeerd.

Dat is bewust: limieten en asymptotisch gedrag moeten uiteindelijk
niet alleen berekend, maar ook geïnterpreteerd en geanalyseerd
kunnen worden.
\end{opmerking}


% ============================================================
% 22 THEMA AFSLUITEN
% ============================================================

\FVDsection{De gereedschapskist is opnieuw klaar}

In dit eerste thema hebben we geen volledig nieuwe theorie opgebouwd.

We hebben het wiskundige gereedschap uit het vijfde jaar opnieuw
scherpgesteld.

We begonnen met gemiddelde verandering:

\[
\frac{\Delta y}{\Delta x}.
\]

Door het interval steeds kleiner te maken kwamen we bij het afgeleid
getal:

\[
f'(a).
\]

Daaruit bouwden we de afgeleide functie:

\[
f'(x),
\]

en vervolgens de tweede afgeleide:

\[
f''(x).
\]

Met die functies konden we het verloop van $f$ analyseren.

En met limieten beschreven we wat er gebeurt nabij bijzondere
$x$-waarden en op oneindig.

De belangrijkste ideeën van Thema 1 vormen zo één geheel:

\[
\boxed{
\text{functie}
\;\longrightarrow\;
\text{verandering}
\;\longrightarrow\;
\text{afgeleiden}
\;\longrightarrow\;
\text{verloop}
\;\longrightarrow\;
\text{limietgedrag}
}
\]

Vanaf nu kunnen we deze gereedschapskist gebruiken bij nieuwe
functieklassen en nieuwe technieken.

\end{document}